\subsection{单 g-模张量积的分解}

命题 2。令 \(\lambda, \mu \in \mathrm{P}_{++}\)。在 \(\mathcal{R}(\mathfrak{g})\) 中，有

\[
[ \lambda ]. [ \mu ] = \sum_ {\nu \in \mathrm{P} _ {+ +}} m (\lambda , \mu , \nu) [ \nu ]
\]

其中

\[
m (\lambda , \mu , \nu) = \sum_ {w, w ^ {\prime} \in \mathrm{W}} \varepsilon (w w ^ {\prime}) \mathfrak {P} (w (\lambda + \rho) + w ^ {\prime} (\mu + \rho) - (\nu + 2 \rho)).
\]

令 E、F 为最高权分别为 \(\lambda, \mu\) 的有限维单 g-模。令 \(l_{\nu}\) 为 E \(E \otimes F\) F 中最高权为 \(\nu\) 的同构型分支的长度。只需证明

\[
l _ {\nu} = \sum_ {w, w ^ {\prime} \in \mathrm{W}} \varepsilon (w w ^ {\prime}) \mathfrak {P} (w (\lambda + \rho) + w ^ {\prime} (\mu + \rho) - (\nu + 2 \rho)).\tag{2}
\]

置 \(c_{1} = \operatorname{ch}(\mathrm{E}) = \sum_{\sigma \in \mathrm{P}} m_{\sigma} e^{\sigma}\)，\(c_{2} = \operatorname{ch}(\mathrm{F})\)，以及 \(d = \mathrm{J}(e^{\rho})\)，其中 \(\mathrm{J}\) 如第2号所定义。有

\[
\sum_ {\xi \in \mathrm{P} _ {+ +}} l _ {\xi} \mathrm{ch} [ \xi ] = \mathrm{ch} (\mathrm{E} \otimes \mathrm{F}) = c _ {1} c _ {2}
\]

因此，在乘以 \(d\) 并应用定理1 后，

\[
\begin{array}{c} \sum_ {\xi \in \mathrm{P} _ {+ +}} l _ {\xi} \mathrm{J} (e ^ {\xi + \rho}) = c _ {1} \mathrm{J} (e ^ {\mu + \rho}) = \left(\sum_ {\sigma \in \mathrm{P}} m _ {\sigma} e ^ {\sigma}\right) \left(\sum_ {w \in \mathrm{W}} \varepsilon (w) e ^ {w (\mu + \rho)}\right) \\ = \sum_ {\tau \in \mathrm{P}} \left(\sum_ {w \in \mathrm{W}} \varepsilon (w) m _ {\tau + \rho - w (\mu + \rho)}\right) e ^ {\tau + \rho}. \end{array}\tag{3}
\]

现在，若 \(\xi \in \mathrm{P}_{++}\)，则 \(\xi + \rho\) 属于由 B 定义的房（第 VI 章第1节第10号）；因此，对所有异于 1 的 \(w \in \mathrm{W}\)，都有 \(w(\xi + \rho) \notin \mathrm{P}_{++}\)。所以，\(e^{\nu + \rho}\) 中 \(\sum_{\xi \in \mathrm{P}_{++}} l_{\xi} \mathrm{J}(e^{\xi + \rho})\) 的系数等于 \(l_{\nu}\)。由 (3) 得

\[
l _ {\nu} = \sum_ {w \in \mathrm{W}} \varepsilon (w) m _ {\nu + \rho - w (\mu + \rho)},
\]

也就是说，由命题1，

\[
l _ {\nu} = \sum_ {w, w ^ {\prime} \in \mathrm{W}} \varepsilon (w) \varepsilon (w ^ {\prime}) \mathfrak {P} (w ^ {\prime} (\lambda + \rho) - (\nu + \rho - w (\mu + \rho) + \rho))
\]

这就证明了 (2)。

例。回到第1号的例子。令 \(\lambda = (n/2)\alpha, \mu = (p/2)\alpha\)，\(\nu = (q/2)\alpha\)，其中 \(n \geq p\)。有

\[
\begin{array}{l} m (\lambda , \mu , \nu) = \mathfrak {P} \left(\frac {n}{2} \alpha + \frac {\alpha}{2} + \frac {p}{2} \alpha + \frac {\alpha}{2} - \frac {q}{2} \alpha - \alpha\right) \\ \qquad - \mathfrak {P} \left(\frac {n}{2} \alpha + \frac {\alpha}{2} - \frac {p}{2} \alpha - \frac {\alpha}{2} - \frac {q}{2} \alpha - \alpha\right) \\ \qquad - \mathfrak {P} \left(- \frac {n}{2} \alpha - \frac {\alpha}{2} + \frac {p}{2} \alpha + \frac {\alpha}{2} - \frac {q}{2} \alpha - \alpha\right) \\ \qquad + \mathfrak {P} \left(- \frac {n}{2} \alpha - \frac {\alpha}{2} - \frac {p}{2} \alpha - \frac {\alpha}{2} - \frac {q}{2} \alpha - \alpha\right) \\ \qquad = \mathfrak {P} \left(\frac {n + p - q}{2} \alpha\right) - \mathfrak {P} \left(\frac {n - p - q - 2}{2} \alpha\right). \end{array}
\]

当 \(n + p + q\) 不被 2 整除，或 \(q \geq n + p\) 时，该数为零。若

\[
q = n + p - 2 r
\]

其中 \(r\) 为 \(\geq 0\) 的整数，则有

\[
m (\lambda , \mu , \nu) = \mathfrak {P} (r \alpha) - \mathfrak {P} ((r - p - 1) \alpha)
\]

因而，当 \(m(\lambda,\mu,\nu)=1\) 时 \(r\leq p\)，当 \(m(\lambda,\mu,\nu)=0\) 时 \textgreater。最后，g-模 \(\mathrm{V}(n)\otimes\mathrm{V}(p)\) 同构于

\[
\mathrm{V} (n + p) \oplus \mathrm{V} (n + p - 2) \oplus \mathrm{V} (n + p - 4) \oplus \dots \oplus \mathrm{V} (n - p)
\]

（Clebsch–Gordan 公式）。

